\begin{proofbox}
	\textbf{\textcolor{CProof}{思路提示}}
	由于点 $D$ 在 $AB$ 上，$E$ 在 $AC$ 上，$\angle A$ 是公共角。在相似三角形中，对应角相等，因此 $\angle A$ 的对应角有三种可能。我们分情况讨论，并证明比例式总归结为两种。

	\medskip
	\textbf{\textcolor{CProof}{正式推导}}
	\begin{description}[style=nextline, leftmargin=2.8em, labelsep=0.5em, itemsep=0.55em]
		\item[\textbf{步骤一（对应角分析）：}]
			在 $\triangle ADE$ 和 $\triangle ABC$ 中，$\angle DAE = \angle BAC$（公共角）。由相似知对应角相等，故 $\angle DAE$ 的对应角有三种可能：$\angle BAC$、$\angle ABC$ 或 $\angle ACB$。
			\par\textbf{情况一：} $\angle DAE$ 对应 $\angle BAC$（即 $A\leftrightarrow A$）。这是最常见的情形，且当 $\angle A$ 不同于 $\angle B,\angle C$ 时唯一可能。
			\par\textbf{情况二：} $\angle DAE$ 对应 $\angle ABC$（即 $A\leftrightarrow B$），此时需 $\angle BAC = \angle ABC$，可得 $AC = BC$。
			\par\textbf{情况三：} $\angle DAE$ 对应 $\angle ACB$（即 $A\leftrightarrow C$），类似可得 $AB = BC$。 \par 下面详细推导情况一，并指出情况二、三可化归为同一比例式。
		\item[\textbf{步骤二（在情况一下分类点 $D$ 的对应）：}]
			由于 $A$ 已固定，且 $D$ 在边 $AB$ 上，$AD$ 与 $AB$ 共线，所以点 $D$ 可能对应 $B$ 或 $C$。
			\par\textbf{理由：} 如果 $D$ 对应 $B$，则 $E$ 自然对应 $C$；如果 $D$ 对应 $C$，则 $E$ 对应 $B$。这是顶点对应的全部可能。
		\item[\textbf{步骤三（列出比例式）：}]
			情形一（$D\leftrightarrow B,\;E\leftrightarrow C$）：
			\[
				\frac{AD}{AB} = \frac{AE}{AC}
			\]
			情形二（$D\leftrightarrow C,\;E\leftrightarrow B$）：
			\[
				\frac{AD}{AC} = \frac{AE}{AB}
			\]
			\par\textbf{理由：} 相似三角形对应边成比例。在情形一中，$AD$ 与 $AB$ 对应，$AE$ 与 $AC$ 对应；在情形二中，$AD$ 与 $AC$ 对应，$AE$ 与 $AB$ 对应。
	\end{description}

	\medskip
	\textbf{\textcolor{CProof}{结论回扣}}
	\textbf{在情况一下我们得到两个比例式。对于情况二、三，利用所得边等关系可验比例式仍化为这两种。因此，在题设条件下，无论哪种对应，$\triangle ADE$ 与 $\triangle ABC$ 相似必满足上述两个比例式之一。解题时需根据具体数据判断哪一个成立或两个均可能。}
\end{proofbox}
